\begingroup
\setlength{\parindent}{0pt}
\setlength{\parskip}{4pt}
\colorlet{panelAccent}{spectreAccent}
\clearpage
\subsection{Input}

Patch the signal you want to follow over time. Spectre combines all voices
of a polyphonic cable into one signal; it does not draw each voice separately.
Voices can reinforce or cancel one another.

\textbf{Input Gain} starts at \textbf{+6 dB}, about twice the amplitude.
It ranges from silence to +12 dB and changes newly captured audio only.
Set it to 0 dB when matching Fourier's default gain. It does not change your
listening path. To adjust the picture instead, use the color range below.

\subsection{Vertical Color Range}

\textbf{Use the color range to decide which levels stand out.} Click the
palette name to choose the colors. Click the scale below it to choose
\textbf{Decibels} or \textbf{Linear}. Decibels makes quiet details easier
to see alongside loud ones; Linear gives strong amplitudes more prominence.

\begin{figure}[!htp]
\centering
% Reuse the panel control geometry, with readable in-context annotations.
\begin{tikzpicture}[x=1pt,y=-1pt,font=\small,line join=round]
  \begin{scope}[shift={(0,-106)}]
    \PanelColorControl
  \end{scope}
  \draw[panel leader] (65,14)--(106,14);
  \node[anchor=west,text=panelInk] at (112,14) {Palette: choose the colors};
  \draw[panel leader] (65,32)--(106,48);
  \node[anchor=west,text=panelInk] at (112,48) {Scale: Decibels or Linear};
  \draw[panel leader] (61,84)--(106,84);
  \node[anchor=west,align=left,text=panelInk] at (112,91)
    {Ceiling: upper-right handle\\Levels above this share the top color};
  \draw[panel leader] (15,147)--(-4,147)--(-4,206)--(106,206);
  \node[anchor=west,align=left,text=panelInk] at (112,199)
    {Floor: lower-left handle\\Levels below this share the bottom color};
\end{tikzpicture}

\caption{The color controls enlarged from the panel reference. Floor and
Ceiling select the levels assigned to the two ends of the palette. The
ramp here is schematic; the chosen palette determines the actual colors.}
\end{figure}

Lower \textbf{Floor} to reveal quiet tails; raise it to hide faint background
detail. Raise \textbf{Ceiling} to distinguish strong peaks that share the top
color. Lower it to devote more contrast to quieter levels. Values outside
the range share the endpoint colors; the stored magnitudes are unchanged.

\paragraph{Try it}
Freeze a sound with Run, then move Floor. The picture changes while the
recorded sound stays the same. More visible background in Decibels mode
usually means you can now see weak content, including noise or leakage;
the color controls do not add noise or apply a noise gate.

\clearpage
\subsection*{Color Range: Everyday Adjustments}

\begin{center}
\renewcommand{\arraystretch}{1.4}
\begin{tabularx}{\textwidth}{@{}p{.38\textwidth}X@{}}
\toprule
\textbf{To do this} & \textbf{Use this gesture} \\
\midrule
Set Floor / Ceiling & Drag the lower-left / upper-right handle. \\
Make a fine adjustment & Hold Cmd (macOS) or Ctrl (other platforms) while dragging. \\
Enter an exact level & Right-click the corresponding handle. \\
Reset one endpoint & Double-click its handle. \\
Undo a drag & Use Rack's Undo; each drag is one action. \\
\bottomrule
\end{tabularx}
\end{center}

Each scale remembers its own range. Decibels starts at \textbf{$-90$ to 0 dB};
Linear starts at \textbf{0--100\%}. The handles cannot cross. A reset stops
at the other endpoint if necessary. Full ranges and palettes are listed on
page~\pageref{sec:control-values}.

Changing the range, scale, or palette recolors existing history, even while
frozen. It neither advances the scan line nor changes the Raw readout.
The \hyperref[sec:detail]{Finding Quiet Detail} exercise walks through a
contrast adjustment on the same captured sound.

\subsection{Run}

\textbf{Lit: record new spectra. Dark: freeze the history.} Press Run to
inspect a moment before the scan line overwrites it. Press again to resume
from the paused position. Run starts enabled; sound received during a pause
is not recorded.

While frozen, change the color controls, frequency view, or Slope to inspect
the same history differently. Gain, AC coupling, Window, Average, and Smooth
need new audio or analysis after you resume. They do not rebuild old columns.
Unlike Fourier, Spectre pauses its analysis as well as its capture.

\paragraph{Keep a result}
Saving a Rack patch saves the settings and Run state, not the image history.
Save an external image if you need a lasting record; a frozen patch reopens
without the old sound and must capture again.

\clearpage
\subsection{Window Function}
A pure tone does not always appear as one thin line. The analyzer works on
short pieces of audio, and the edges of each piece can spread the tone into
nearby frequencies. This spreading is called \textit{leakage}. A window
softens the edges, trading a broader main peak for less energy around it.

\textbf{Start with Flattop} when comparing the heights of steady tones.
Try \textbf{Hann} when you want to distinguish nearby harmonics. Boxcar gives
a narrow main peak, but its leakage can hide quieter tones beside a loud one.
No window is best for every sound. The FFT measures at evenly spaced
frequencies called \textit{bins}; the plots below show what it sees.

\begin{figure}[!htp]
\centering
\input{WindowTradeoffs.tex}
\caption{Calculated responses to one tone halfway between FFT bins.
The dashed line marks the tone; dots mark measured bins. All three plots use
the same 2048-sample length and level scale. The ripples are leakage, not
additional tones.}
\end{figure}

\paragraph{Read the shape, not just the height}
The small ripples around a large peak can come from the analysis itself;
they are not necessarily extra notes or noise in the source. Compare windows
before drawing that conclusion. The module compensates for the window's
level reduction, but peak shape and the displayed noise level still change.

\paragraph{Try it}
With Average and Smooth off, capture a steady oscillator using Hann, then
switch to Flattop. Compare the width of the new band with the older history.
Window changes apply to \textbf{new columns only}; they do not redraw the
past. The complete menu is on page~\pageref{sec:control-values}.

\clearpage
\subsection{Freq Scale}

Choose \textbf{Linear} to compare harmonic spacing: equal steps in hertz
occupy equal distances from bottom to top. Choose \textbf{Logarithmic}
(the default) to give bass frequencies more room. Frequency is vertical in
Spectre; the colors show level.

The control's Logarithmic mode uses a square-root mapping in this version,
so octaves do not occupy equal heights. Read the labels or hover readout
when you need an exact frequency.

\subsection{Average}
Use \textbf{Average} to calm a flickering display. Increasing it makes
changes settle more gradually, which helps compare the overall balance of
noise or a busy mix. Leave it at \textbf{0 ms} when following attacks and
short events. The range extends to 2500 ms.

The value controls how strongly previous results influence the next result;
it does not select a fixed-length recording to average.

In Spectre, Average can blur an attack across successive columns. It does
not change the duration of the history, and it affects new columns only.

\subsection{Smooth}
Use \textbf{Smooth} to see broad tonal balance instead of every small
peak. Start at \textbf{None} for individual harmonics. Try 1/12 octave for a
gentler outline, then increase the width if you want less detail. Wider
settings can lower narrow peaks and merge neighboring ones.

An octave means a doubling of frequency. An octave-based smoothing width
therefore covers more hertz in the treble than in the bass. Smooth does not
create additional frequency resolution.

\begin{figure}[!htp]
\centering
\input{Smoothing.tex}
\caption{Average softens changes across time; Smooth blends detail across
frequency. These are qualitative sketches, not measured responses.}
\end{figure}

Like Average, Smooth affects new columns only. Allow a fresh sweep
before comparing settings.

\clearpage
\subsection{Low Frequency}

\textbf{LO Freq} sets the bottom of the frequency view. Raise it to leave
frequencies below your area of interest out of the picture. It starts at
\textbf{0 Hz}. It does not filter the input or erase stored frequencies.

\subsection{High Frequency}

\textbf{HI Freq} sets the top of the view. Lower it to give the bass or a
group of harmonics more room. Keep it above LO Freq. Both controls can reach
half of Rack's sample rate, called the \textit{Nyquist frequency}; that is
24 kHz at a 48 kHz sample rate. HI Freq starts at this upper limit.

\paragraph{Try it}
Freeze a sound and view only 0--500 Hz. The bands grow taller on screen,
but the underlying frequency detail is unchanged. Spectre has a fixed
analysis length: if two bass notes merge, use Fourier with a longer Length
to separate them. Check the view after a sample-rate change.

\subsection{Slope}

\textbf{Slope changes which frequencies stand out in color.} Positive
values emphasize high frequencies; negative values emphasize low ones.
It can help reveal quiet upper harmonics. Set it to \textbf{0} for a fair
comparison of unweighted levels.

The default is \textbf{+4.5 dB per octave}, and the range is $-9$ to
$+9$ dB per octave. An octave doubles the frequency. In Decibels mode,
1 kHz is the pivot: for example, +3 dB/oct adds 3 dB at 2 kHz and subtracts
3 dB at 500 Hz. The DC bin receives no tilt.

Slope recolors retained history, including while frozen. It changes the
\textbf{Color} readout, but leaves stored magnitudes and \textbf{Raw} alone.
It never changes the audio in your listening path.

\paragraph{A note about Linear intensity}
Linear intensity keeps the earlier display behavior. With Linear Freq Scale,
Slope pivots around 1 kHz. With Logarithmic Freq Scale, it follows image-row
position instead of exact frequency, so the tilt does not match the stated
dB per octave at the frequency labels. Use Decibels intensity for frequency-based
tilt, or zero Slope for an unweighted comparison.

\endgroup
